\section{Limitations}
\label{sec:limits}

Two smooth Hamiltonian systems --- a pendulum and a two-degree-of-freedom anharmonic oscillator ---
with no contacts and no multiple interacting objects; whether any of this survives discontinuous
dynamics is untested. Three trained models per arm, with the model seed as the independent unit and
ranges rather than confidence intervals, because three is too few for a precise effect size. The
extraction minimises an invariance ratio that \emph{is} the normal component of one-step error, so a
low normal error is the objective being met, not a finding; the empirical content is that the
resulting direction coincides with physical energy, which supervision does not reproduce. The
architecture comparison likewise does not isolate architecture: an open-loop reconstruction term and
a prior--posterior KL are different losses, not the same loss at different strengths.

\paragraph{No downstream benefit, and we measured why.}
Both practical uses we tested fail, for one reason. Planning with CEM over the learned model on an
energy-targeting task, the model plans well --- beating random actions by $+1.65$, 95\% CI
$[+0.50,+2.79]$ --- yet no correction arm differs from no correction over three models and 20
episodes each: the largest effect is $0.0014$ against a return standard deviation of $0.575$,
\textbf{0.24\% of one episode's spread}, with all four registered predictions failing. The cause is
structural. The correction acts on \emph{accumulated} drift, and at the planner's horizon it shifts
imagined energy by $0.3\%$ of the spread across candidate action sequences --- the quantity the
planner ranks plans by --- so it cannot change which plan is chosen. Our repair is measured over 100
autonomous steps; a re-planning controller never imagines long enough to accumulate what the method
removes, and looking further ahead does not help, because planning quality itself degrades with
horizon. The same effect-size limit defeats drift as an online trust signal, which is beaten by plain
latent displacement on every model. Appendix~\ref{sec:causal-detail} gives both in full.

\paragraph{The method finds only what the data gave the model a reason to represent.}
The pendulum's rod length is fixed, an exact constraint, yet searching for a residual $G(z)=0$ fails:
the recovered direction correlates with the true residual at $0.004$--$0.032$. Minimising
$\mathbb{E}[G^2]$ is the wrong objective, ranking the true constraint $1525$th of $1819$ directions,
and the constraint is barely encoded at all --- predicted from the latent at held-out $|\rho|=0.19$.
\textbf{A quantity that never varies carries no information}, so nothing in a reconstruction
objective pushes it into the latent; it can live in the decoder's weights. This bounds the approach
to quantities that \emph{vary}, and is why we report no released-checkpoint experiment: limb lengths
do not vary either.

Finally, the two-invariant case limits the extraction rather than the model. Where both energy and
angular momentum are conserved, the conserved subspace reliably contains both (angular-momentum
correlation $0.94$--$0.98$) while the joint fit isolates either in only 1 of 7 measurements: flow
alignment selects a single direction inside that subspace and has no unique answer when two exist.
